\documentclass[11pt,a4paper]{article}
\usepackage[utf8]{inputenc}
\usepackage[T1]{fontenc}
\usepackage[ngerman]{babel}
\usepackage{geometry}
\geometry{margin=2.5cm}
\usepackage{booktabs}
\usepackage{enumitem}
\usepackage{xcolor}
\usepackage{scrextend}
\usepackage{hyperref}
\usepackage{tabularx}
\hypersetup{
	colorlinks=true,
	linkcolor=blue,
	citecolor=blue,
	urlcolor=blue,
}

\title{Alkohol vs. Cannabis: Ein wissenschaftlicher Risikovergleich}
\author{Zusammenstellung auf Basis aktueller Forschung}
\date{\today}

\begin{document}
	
	\maketitle
	
	\begin{abstract}
		Die Frage nach der Gefährlichkeit von Alkohol versus Cannabis ist komplex, da beide Substanzen auf unterschiedliche Weise Schaden anrichten können. Aus einer medizinischen und gesellschaftlichen Gesamtperspektive wird Alkohol jedoch von der überwiegenden Mehrheit der wissenschaftlichen Studien als die \textbf{deutlich gefährlichere Droge} eingestuft. Dieses Dokument fasst die zentralen Erkenntnisse zusammen.
	\end{abstract}
	
	\section{Einleitung}
	
	Während Alkohol in vielen Gesellschaften tief verwurzelt ist, gilt Cannabis oft als relativ harmlos. Ein direkter Vergleich der Risiken zeigt jedoch ein klares Bild: Alkohol ist in nahezu jeder Kategorie gefährlicher – für das Individuum und für die Gesellschaft.
	
	\section{Akute Toxizität und Überdosierung}
	
	\begin{itemize}[leftmargin=*, noitemsep]
		\item \textbf{Alkohol:} Eine akute Alkoholvergiftung kann tödlich enden. Weltweit und auch in Deutschland (ca. 44.000 Todesfälle jährlich\cite{who-deaths}\cite{dhs-deaths}) sind Todesfälle durch Alkoholkonsum eine traurige Realität. Schätzungen der WHO zufolge wird in der Europäischen Region jeder elfte Todesfall durch Alkohol verursacht\cite{who-deaths}.
		\item \textbf{Cannabis:} Eine tödliche Überdosis ist extrem unwahrscheinlich. Das \textit{National Institute on Drug Abuse} (NIDA) bestätigt, dass bisher kein Todesfall ausschließlich auf eine Cannabis-Überdosis zurückzuführen ist\cite{nida-no-overdose}. Auch in Deutschland sind keine entsprechenden Fälle verzeichnet\cite{bmg-cannabis}.
	\end{itemize}
	
	\section{Langfristige körperliche Schäden}
	
	\begin{itemize}[leftmargin=*, noitemsep]
		\item \textbf{Alkohol:} Alkohol ist ein Zellgift und eine Hauptursache für Leberzirrhose, Pankreatitis, Herzmuskelerkrankungen und verschiedene Krebsarten. Die WHO stuft Alkohol als \textbf{Gruppe-1-Karzinogen} ein (höchste Risikogruppe, gemeinsam mit Asbest und Tabak)\cite{who-carcinogen}. Er ist kausal mit mindestens sieben Krebsarten verbunden, darunter Mundhöhlen-, Rachen-, Speiseröhren-, Leber-, Kehlkopf-, Darm- und Brustkrebs\cite{who-carcinogen}.
		\item \textbf{Cannabis:} Die direkten körperlichen Schäden sind weniger schwerwiegend. Das größte Risiko für chronische Schäden besteht durch das \textbf{Rauchen} (oft mit Tabak gemischt), was zu Atemwegserkrankungen führen kann. Es gibt Hinweise auf ein erhöhtes Risiko für Herz-Kreislauf-Erkrankungen\cite{nida-heart}. Ein eindeutiger Nachweis für ein Krebsrisiko durch Cannabis allein liegt nicht vor.
	\end{itemize}
	
	\section{Psychische Gesundheit und Gehirn}
	
	\begin{itemize}[leftmargin=*, noitemsep]
		\item \textbf{Alkohol:} Chronischer Konsum kann zu alkoholbedingter Demenz, Gedächtnisstörungen, Depressionen und Angststörungen führen. Ein körperlicher Alkoholentzug kann \textbf{lebensbedrohlich} sein (Delirium tremens).
		\item \textbf{Cannabis:} Das zentrale Risiko ist ein \textbf{erhöhtes Psychose-Risiko}, besonders bei Jugendlichen und jungen Erwachsenen mit entsprechender Veranlagung. Eine kanadische Studie zeigte, dass Jugendliche ein 11-fach höheres Risiko für eine psychotische Störung haben, wenn sie Cannabis konsumieren\cite{cannabis-psychosis}. Hochpotentes Cannabis verdoppelt das Risiko für psychotische Symptome im Vergleich zu niedrigpotentem Cannabis zusätzlich\cite{cannabis-high-potency}. Zudem kann Cannabis die Gehirnentwicklung bei Jugendlichen beeinträchtigen.
	\end{itemize}
	
	\section{Abhängigkeitspotenzial}
	
	\begin{itemize}[leftmargin=*, noitemsep]
		\item \textbf{Alkohol:} Hat ein hohes körperliches und psychisches Abhängigkeitspotenzial. In Deutschland sind schätzungsweise 1,6 Millionen Menschen alkoholabhängig\cite{dhs-abhaengigkeit}.
		\item \textbf{Cannabis:} Führt vor allem zu einer psychischen Abhängigkeit, körperliche Entzugserscheinungen sind milder. Etwa \textbf{9\% aller Konsumenten} entwickeln eine Abhängigkeit, wobei die Rate bei jugendlichen Konsumenten auf 17\% steigt\cite{dhs-cannabis-abhaengigkeit}.
	\end{itemize}
	
	\section{Gesellschaftliche und indirekte Schäden}
	
	\begin{itemize}[leftmargin=*, noitemsep]
		\item \textbf{Alkohol:} Verursacht immense gesellschaftliche Kosten durch Unfälle (ca. 30\% aller tödlichen Verkehrsunfälle), Gewalt, Kriminalität und Produktivitätsausfall. Die Deutschen Hauptstelle für Suchtfragen (DHS) beziffert die volkswirtschaftlichen Kosten des Alkoholkonsums in Deutschland auf rund \textbf{57 Milliarden Euro pro Jahr}\cite{dhs-costs}. In umfassenden Schadensskalen erreicht Alkohol Spitzenwerte, Cannabis wird deutlich niedriger bewertet.
		\item \textbf{Cannabis:} Die gesellschaftlichen Schäden sind im Vergleich zu Alkohol geringer. Die volkswirtschaftlichen Kosten sind schwer zu beziffern, liegen aber deutlich unter denen von Alkohol.
	\end{itemize}
	
	\section{Zusammenfassung der Risikoprofile}
	
	\begin{table}[h!]
		\centering
		\small                     % Optional: Schrift etwas verkleinern für mehr Platz
		\begin{tabularx}{\textwidth}{@{} l X X @{}}
			\toprule
			\textbf{Kategorie} & \textbf{Alkohol} & \textbf{Cannabis} \\ \midrule
			Akute Toxizität  & Hoch, Überdosis kann tödlich sein & Sehr gering, tödliche Überdosis extrem unwahrscheinlich\cite{nida-no-overdose} \\
			Organschäden      & Schwer (Leber, Herz, Gehirn, Pankreas) & Geringer, Hauptrisiko durch Rauchen (Lunge) \\
			Krebsrisiko       & Eindeutig krebserregend (WHO Gruppe 1)\cite{who-carcinogen} & Kein eindeutiger Nachweis, Risiko durch Tabak-Mischkonsum \\
			Psychische Risiken & Demenz, Depression, lebensgefährlicher Entzug & Erhöhtes Psychoserisiko, besonders bei Jugendlichen\cite{cannabis-psychosis} \\
			Abhängigkeit      & Hoch, körperlich und psychisch\cite{dhs-abhaengigkeit} & Mittel, vorwiegend psychisch (9\%-17\%)\cite{dhs-cannabis-abhaengigkeit} \\
			Gesellschaftsschaden & Sehr hoch (57 Mrd. € jährlich in D)\cite{dhs-costs} & Geringer im Vergleich zu Alkohol \\ \bottomrule
		\end{tabularx}
	\end{table}
	
	\section{Fazit}
	
	Die wissenschaftliche Evidenz zeigt, dass Alkohol in den meisten medizinisch relevanten Kategorien das deutlich höhere Risikoprofil aufweist. Dennoch ist Cannabis keine harmlose Substanz – beide Drogen haben das Potenzial, der Gesundheit zu schaden, und der sicherste Konsum ist immer der Verzicht.
	
	\newpage
	\begin{thebibliography}{99}
		\bibitem{who-deaths} Weltgesundheitsorganisation (WHO), 2025: \href{https://www.who.int/europe/de/news/item/24-12-2025-1-in-3-injury-deaths-caused-by-alcohol--who-europe-reiterates-the-harms}{WHO/Europa weist erneut auf die schädlichen Folgen hin}.
		\bibitem{dhs-deaths} Deutsche Hauptstelle für Suchtfragen (DHS), \textit{Jahrbuch Sucht 2026}, Kapitel 2.1.
		\bibitem{nida-no-overdose} National Institute on Drug Abuse (NIDA); zitiert nach \href{https://www.dw.com/de/faktencheck-wie-ungef%C3%A4hrlich-ist-cannabis/a-65853237}{Deutsche Welle (DW), Faktencheck: Wie (un)gefährlich ist Cannabis?}
		\bibitem{bmg-cannabis} Bundesministerium für Gesundheit: \href{https://www.bundesgesundheitsministerium.de/themen/cannabis/faq-cannabis/ueberdosierung.html}{Cannabis: Fragen und Antworten zur Überdosierung}.
		\bibitem{who-carcinogen} Weltgesundheitsorganisation (WHO), 2021: \href{https://www.who.int/europe/de/home/20-10-2021-alcohol-is-one-of-the-biggest-risk-factors-for-breast-cancer}{Alkohol ist einer der größten Risikofaktoren für Brustkrebs}; IARC-Klassifizierung von 1988.
		\bibitem{nida-heart} National Institute on Drug Abuse (NIDA): \href{https://nida.nih.gov/publications/research-reports/marijuana/cardiovascular-health-effects}{Cardiovascular Health Effects of Marijuana}.
		\bibitem{cannabis-psychosis} André McDonald et al., 2025: Studie zum 11-fach erhöhten Psychoserisiko bei jugendlichen Cannabiskonsumenten, zitiert nach \href{https://www.drugcom.de/news/jugendliche-die-cannabis-konsumieren-haben-ein-erhoehtes-risiko-fuer-psychose/}{drugcom.de}.
		\bibitem{cannabis-high-potency} Lindsey Hines et al., 2024: Längsschnittstudie zu hochpotentem Cannabis und Psychoserisiko, zitiert nach \href{https://www.drugcom.de/news/hochpotenter-cannabis-verdoppelt-risiko-fuer-psychotische-symptome/}{drugcom.de}.
		\bibitem{dhs-abhaengigkeit} Deutsche Hauptstelle für Suchtfragen (DHS): \href{https://www.dhs.de/sucht/zahlen-und-fakten.html}{Zahlen und Fakten zu Alkoholabhängigkeit}.
		\bibitem{dhs-cannabis-abhaengigkeit} Deutsche Hauptstelle für Suchtfragen (DHS), \textit{Jahrbuch Sucht 2025}, Kapitel 2.5; Techniker Krankenkasse: \href{https://www.tk.de/techniker/gesundheit-und-medizin/erste-hilfe-und-gesundheitsvorsorge/sucht-und-abhaengigkeit/cannabis/cannabis-abhaengigkeitspotenzial-2075314}{Cannabis-Abhängigkeitspotenzial}.
		\bibitem{dhs-costs} Deutsche Hauptstelle für Suchtfragen (DHS): \href{https://www.dhs.de/sucht/zahlen-und-fakten.html}{Volkswirtschaftliche Kosten des Alkoholkonsums}.
	\end{thebibliography}
\end{document}